% What the period handler does now, and the one path that does not pass
% through it.
\begin{tikzpicture}[font=\sffamily\scriptsize, x=1cm, y=1cm]

\foreach \x/\name/\id in {0/{the decoder}/dec, 3.3/{period handler}/isr,
                          6.6/{the timer}/tim, 9.9/{plant model}/pl,
                          13.2/{generator}/gen}{
  \node[actor, text width=21mm] (\id) at (\x,0) {\name};
  \draw[lifeline] (\id.south) -- ++(0,-9.3);
}

\draw[msg] (0,-1.2) -- node[above]{position, through the real decode path} (3.3,-1.2);
\node[activation, minimum height=5.6mm] at (3.3,-1.5) {};

\draw[msg] (3.3,-2.2) -- ++(1.0,0) -- ++(0,-0.5)
  node[right, xshift=1mm]{controller: empty until chapter 16} -- ++(-1.0,0);

\draw[msg] (3.3,-3.6) -- node[above]{duty, into the compare registers} (6.6,-3.6);
\node[activation, minimum height=3mm] at (6.6,-3.8) {};
\node[tlabel, anchor=west] at (6.75,-4.3) {dead time inserted in hardware};

\draw[msg] (3.3,-4.9) -- node[above]{step once, dt = the control period} (9.9,-4.9);
\node[activation, minimum height=3.4mm] at (9.9,-5.15) {};
\draw[reply] (9.9,-5.7) -- node[below]{modelled velocity} (3.3,-5.7);

\draw[msg] (3.3,-6.5) -- node[above]{set the rate and the direction} (13.2,-6.5);
\node[activation, minimum height=3mm] at (13.2,-6.7) {};
\node[tlabel, anchor=east] at (13.05,-7.2) {three wires, out of the chip};

% the hardware-only path, drawn to the right so its label clears the lifelines
\draw[faultedge] (6.6,-7.9) -- ++(1.1,0) -- ++(0,-0.5)
  node[right, xshift=1mm, text=red!60!black, align=left]
  {break input: the outputs go inactive here,\\with no software involved at all} -- ++(-1.1,0);
\draw[reply, draw=red!70!black] (6.6,-8.9) -- node[below, text=red!60!black]{the handler finds out afterwards} (3.3,-8.9);

\node[umlnote, text width=54mm, anchor=north west] at (0,-9.8)
  {The break path is hardware and software finds out second. A design in which
   software must run for the drive to stop is a design whose safety depends on
   the scheduler, and chapter 18 formalises that ordering for every other fault
   the node can have.};
\node[umlnote, text width=54mm, anchor=north west] at (7.4,-9.8)
  {Nothing reads the plant's state directly. Its position reaches the loop only
   through the generator and the decoder. The single exception is the rig in
   chapter 20, which reads both on purpose and reports the difference.};
\end{tikzpicture}
